
\subsubsection{2.1 Benchmark Saturation and Overfitting Measurement}

Roelofs et al. (2019) quantified benchmark overfitting in CV through retesting studies, demonstrating that CIFAR-10 and ImageNet models show systematic accuracy drops on newly constructed test sets consistent with implicit overfitting to statistical properties of the original test distributions. Recht et al. (2019) independently confirmed this robustness gap on ImageNet. Both works are domain-specific (CV) and retrospective, requiring held-out test sets unavailable in standard leaderboard settings.

Polo et al. (2026) introduced the S-index (S_index = exp(-R_norm²)), a scalar saturation measure for LLM benchmarks that achieves Bayesian R² = 0.884 in predicting saturation from benchmark metadata features. This work is LLM-specific and does not validate a temporal mechanism linking submission accumulation to compression across multiple domains.

\subsubsection{2.2 Test Set Contamination}

Chan et al. (2024) investigated pre-training contamination impact on 75 Kaggle benchmarks (MLE-bench), finding that contamination measurably inflates apparent benchmark performance. ConStat (eth-sri) provides performance-based contamination detection for NLP benchmarks without requiring text access. These tools address contamination as a mechanism of inflation but do not measure the broader compression phenomenon captured by the present work.

\subsubsection{2.3 Dataset Documentation and FAIR Principles}

Gebru et al. (2018) established the Datasheets for Datasets standard for structured dataset documentation. Wilkinson et al. (2016; 2024) formalized FAIR (Findable, Accessible, Interoperable, Reusable) principles for scientific data. Phase 1 research for this project identified that FAIR compliance metrics and benchmark health indicators exist in isolation, with no unified pipeline. The present work focuses on the dynamic monitoring aspect of benchmark health from leaderboard submission trajectories rather than static documentation standards.

\subsubsection{2.4 Survival Analysis in Evaluation Research}

Survival analysis methods including the Cox proportional hazards model (Cox, 1972) and Kaplan-Meier estimators (Kaplan and Meier, 1958) have been applied in clinical and reliability domains but not previously to benchmark lifecycle prediction. The BCBHS framework introduced the framing of benchmark health monitoring as a survival analysis problem. As described in this paper, the survival model component was not executed in this work due to a collapse event operationalization failure.

